\documentclass[11pt,a4paper]{article}
\usepackage[margin=1in]{geometry}
\usepackage{amsmath,booktabs,hyperref,float}

\title{Iteration 014: Long FIR + Dropout + Warm Restarts}
\author{Autoresearch Loop}
\date{April 8, 2026}

\begin{document}
\maketitle

\begin{abstract}
Sweeping longer FIR filter lengths (7, 11, 15 taps) with channel dropout and
cosine warm restarts achieves $r=0.375$, not improving over the 7-tap baseline
($r=0.376$). Longer filters do not capture additional useful temporal structure
in the 1--9\,Hz band at 20\,Hz sampling.
\end{abstract}

\section{Hypothesis}
Longer FIR filters (up to 750\,ms at 15 taps) might capture slower EEG dynamics
or longer volume conduction delays. Warm restarts could help escape local minima.

\section{Method}
Swept $L \in \{7, 11, 15\}$ taps with 15\% channel dropout. Used cosine
annealing with warm restarts ($T_0=50$, $T_\text{mult}=1$) over 150 epochs.
Best result selected by validation correlation.

\section{Results}
\begin{table}[H]
\centering
\begin{tabular}{lrrr}
\toprule
Model & Mean $r$ & Std $r$ & SNR (dB) \\
\midrule
FIR + dropout (iter011) & 0.376 & 0.076 & 0.61 \\
Long FIR + restarts & 0.375 & 0.076 & 0.61 \\
\bottomrule
\end{tabular}
\end{table}

\section{Analysis}
At 20\,Hz with 1--9\,Hz content, the Nyquist-limited temporal resolution means
7 taps (350\,ms) already captures most relevant temporal structure. Longer filters
add parameters that overfit without gaining signal. This confirms 7 taps as the
optimal filter length for this sampling rate.

\section{Next}
Iteration 015: correlation-based loss function to directly optimize the evaluation metric.

\end{document}
